% BEGIN REV09_RECTAS27_COORDENADAS
\subsection{Coordenadas reversibles de los pares de rectas}
\label{corpus9:xii:rectas-coordenadas}

La tabla~\ref{p12-83-c20-s6:tab:h3-27-rectas-u029} fija la biyección
\(\iota:H_3\to\mathscr L_{27}\), donde \(\mathscr L_{27}\) es la
configuración de las veintisiete rectas. Conservamos la ley central
\eqref{p12-83-c20-s6:eq:heisenberg-3-nuevo}:
\[
 (a,b,z)(c,d,w)=(a+c,b+d,z+w+bc),\qquad
 (a,b,z)^{-1}=(-a,-b,-z+ab).
\]
La fórmula de la inversa se comprueba multiplicando en ambos órdenes;
las tres coordenadas resultantes son cero. Todas las operaciones de este
apartado se efectúan en \(\mathbb F_3\).

\begin{proposition}[Codificación reversible de las 729 palabras]
\label{corpus9:xii:rectas-biyeccion}
Para \(w=(\ell_1,h_1,\ell_2,h_2,\ell_3,h_3)\), sean
\(L(w)=(\ell_1,\ell_2,\ell_3)\) y
\(H(w)=(h_1,h_2,h_3)\). La aplicación
\begin{equation}
 \mathcal E:\mathbb F_3^6\longrightarrow
                 \mathscr L_{27}\times\mathscr L_{27},\qquad
 \mathcal E(w)=(\iota L(w),\iota H(w))
 \label{corpus9:xii:rectas-codificacion}
\end{equation}
es biyectiva. En las coordenadas relativas
\(L=(a,b,z)\), \(R=L^{-1}H=(u,v,t)\), su reconstrucción es
\begin{equation}
 H=(a+u,b+v,z+t+bu),\qquad
 w=(a,a+u,b,b+v,z,z+t+bu).
 \label{corpus9:xii:rectas-reconstruccion}
\end{equation}
\end{proposition}
\begin{proof}
Dado un par ordenado \((\ell,h)\) de rectas, la tabla fija los
únicos triples \(L=\iota^{-1}(\ell)\) y
\(H=\iota^{-1}(h)\). Entrelazar sus coordenadas recupera una única
palabra \(w\), y la evaluación de
\eqref{corpus9:xii:rectas-codificacion} devuelve el par inicial.
Recíprocamente, desentrelazar y volver a entrelazar una palabra conserva
sus seis coordenadas. Esto establece las dos composiciones inversas.

La transformación \((L,H)\mapsto(L,L^{-1}H)\) es biyectiva, con
inversa \((L,R)\mapsto(L,LR)\). La ley central da
\(LR=(a+u,b+v,z+t+bu)\), que prueba
\eqref{corpus9:xii:rectas-reconstruccion}. En particular, el término
\(bu\) conserva el cociclo en la tercera coordenada. Si se lo omite,
la reconstrucción cambia precisamente en \(bu\); por ejemplo,
\(L=(0,1,0)\), \(R=(1,0,0)\) producen
\(H=(1,1,1)\), mientras la suma componente a componente produce
\((1,1,0)\).
\end{proof}

La partición del teorema~\ref{p12-83-c20-s6:thm:particion-hensel-schlafli-nueva}
se conserva bajo esta biyección: el elemento relativo \(R\) distingue
diagonal, incidencia y alabeamiento, con cardinales \(27\), \(270\)
y \(432\). Los pares ordenados constituyen el dominio completo de
la coordenatización. La condición de adyacencia es una condición adicional sobre
los pares o sus sucesiones. La prolongación a historias completas se
construye en el teorema~\ref{corpus9:xii:rectas-homeomorfismo} mediante
el lector de bloques del TPK.
% END REV09_RECTAS27_COORDENADAS
